%document
\documentclass[10pt]{beamer}
%theme
\usetheme{metropolis}
% packages
\usepackage{color}
\usepackage{listings}
\usepackage[ngerman]{babel}
\usepackage[utf8]{inputenc}
\usepackage{multicol}

\usepackage{hyperref}
\hypersetup{colorlinks=true, linkcolor=blue, urlcolor=red}

% color definitions
\definecolor{mygreen}{rgb}{0,0.6,0}
\definecolor{mygray}{rgb}{0.5,0.5,0.5}
\definecolor{mymauve}{rgb}{0.58,0,0.82}

\lstset{
    backgroundcolor=\color{white},
    % choose the background color;
    % you must add \usepackage{color} or \usepackage{xcolor}
    basicstyle=\footnotesize\ttfamily,
    % the size of the fonts that are used for the code
    breakatwhitespace=false,
    % sets if automatic breaks should only happen at whitespace
    breaklines=true,                 % sets automatic line breaking
    captionpos=b,                    % sets the caption-position to bottom
    commentstyle=\color{mygreen},    % comment style
    % deletekeywords={...},
    % if you want to delete keywords from the given language
    extendedchars=true,
    % lets you use non-ASCII characters;
    % for 8-bits encodings only, does not work with UTF-8
    frame=single,                    % adds a frame around the code
    keepspaces=true,
    % keeps spaces in text,
    % useful for keeping indentation of code
    % (possibly needs columns=flexible)
    keywordstyle=\color{blue},       % keyword style
    % morekeywords={*,...},
    % if you want to add more keywords to the set
    numbers=left,
    % where to put the line-numbers; possible values are (none, left, right)
    numbersep=5pt,
    % how far the line-numbers are from the code
    numberstyle=\tiny\color{mygray},
    % the style that is used for the line-numbers
    rulecolor=\color{black},
    % if not set, the frame-color may be changed on line-breaks
    % within not-black text (e.g. comments (green here))
    stepnumber=1,
    % the step between two line-numbers.
    % If it's 1, each line will be numbered
    stringstyle=\color{mymauve},     % string literal style
    tabsize=4,                       % sets default tabsize to 4 spaces
    % show the filename of files included with \lstinputlisting;
    % also try caption instead of title
    language = C,
	showspaces = false,
	showtabs = false,
	showstringspaces = false,
	escapechar = ,
}

\def\ContinueLineNumber{\lstset{firstnumber=last}}
\def\StartLineAt#1{\lstset{firstnumber=#1}}
\let\numberLineAt\StartLineAt



\newcommand{\codeline}[1]{
	\alert{\texttt{#1}}
}

% This Document contains the information about this course.

% Authors of the slides
\author{
    Lecturers:
	%name1
	Mirko Jantschke,
	%name2
	Pascal Scholz
	\\
	%creaters of the slides
	Created by: Richard Mörbitz, Manuel Thieme
}


% Fancy Logo
\titlegraphic{\hfill\includegraphics[height=1.25cm]{../templates/fsr_logo_cropped}}


\title{Dynamic storage}
\subtitle{Variables}
\date{\today}

\usetikzlibrary{tikzmark}
\usetikzlibrary{arrows}
\usetikzlibrary{decorations.pathmorphing}
\tikzset{arrow/.style={-latex, shorten >=.5ex, shorten <=.5ex}}

% the actual document
\begin{document}

\maketitle

%-------------------------------------------------------------------------------------------------
\begin{frame}{Contents}
	\tableofcontents
\end{frame}

%-------------------------------------------------------------------------------------------------
\section{Memory allocation}
%-------------------------------------------------------------------------------------------------

\begin{frame}[fragile]{Runtime conditions}
	Think about a helper function that initializes a large portion of memory (e.g. for an array).\\
	\bigskip
	Until now, we have let the compiler decide how to place the variables and arrays in memory. In many cases, this is not sufficient.\\
	Now we want to allocate memory \textbf{explicitely} and \textbf{dynamically}.\\
	\bigskip
	There are four handy functions declared in \textit{stdlib.h}:
	\begin{itemize}
		\item \textit{malloc()}: Allocate a block of memory
		\item \textit{calloc()}: Allocate a block of memory and initialize it
		\item \textit{realloc()}: Alter the size of a block of memory
		\item \textit{free()}: Release a block of memory
	\end{itemize}
\end{frame}

%-------------------------------------------------------------------------------------------------

\begin{frame}[fragile]{A closer look at memory}
	\begin{tikzpicture}[y=2cm, font=\footnotesize]
		\node[above, font=\small] at (1.5,0) {Stack};
		
		\draw (0,0) -- (0,-2);
		\draw (0,0) -- (3,0);
		\draw (3,0) -- (3,-2);
		\draw[line join=round, decoration={zigzag,segment length=4, amplitude=1}, decorate] (0,-2) -- (3,-2);
		
		\draw[->] (-.25,0) -- (-.25,-1.9);
		
		\draw[dashed] (0,-.25) -- (3,-.25);
		\draw[dashed] (0,-.5) -- (3,-.5);
		\draw[dashed] (0,-.75) -- (3,-.75);
		\draw[dashed] (0,-1) -- (3,-1);
		\draw[dashed] (0,-1.25) -- (3,-1.25);
		\draw[dashed] (0,-1.5) -- (3,-1.5);
		\draw[dashed] (0,-1.75) -- (3,-1.75);
		
		\node[below, font=\small] at (6.5,-2) {Heap};		
		
		\draw (5,0) -- (5,-2);
		\draw[line join=round, decoration={zigzag,segment length=4, amplitude=1}, decorate] (5,0) -- (8,0);
		\draw (8,0) -- (8,-2);
		\draw (5,-2) -- (8,-2);
		
		\draw[<-] (8.25,-.1) -- (8.25,-2);
		
		\draw[dashed] (5,-.25) -- (8,-.25);
		\draw[dashed] (5,-.5) -- (8,-.5);
		\draw[dashed] (5,-.75) -- (8,-.75);
		\draw[dashed] (5,-1) -- (8,-1);
		\draw[dashed] (5,-1.25) -- (8,-1.25);
		\draw[dashed] (5,-1.5) -- (8,-1.5);
		\draw[dashed] (5,-1.75) -- (8,-1.75);
		
		\draw[thick, dashed, orange] (0,0) -- (3,0);
		\draw[thick, dashed, orange] (0,0) -- (0,-.25);
		\draw[thick, dashed, orange] (0,-.25) -- (3,-.25);
		\draw[thick, dashed, orange] (3,0) -- (3,-.25);
		\node[orange,below=.15em] at (1.5,0){int};

		\draw[thick, dashed, orange] (0,-.25) -- (3,-.25);
		\draw[thick, dashed, orange] (0,-.25) -- (0,-.5);
		\draw[thick, dashed, orange] (0,-.5) -- (3,-.5);
		\draw[thick, dashed, orange] (3,-.25) -- (3,-.5);
		\node[orange,below=.15em] at (1.5,-.25){int};
		
		\begin{uncoverenv}<3->
			\draw[thick, dashed, orange] (5,-1.75) -- (8,-1.75);
			\draw[thick, dashed, orange] (5,-1.75) -- (5,-2);
			\draw[thick, dashed, orange] (5,-2) -- (8,-2);
			\draw[thick, dashed, orange] (8,-1.75) -- (8,-2);
			\node[orange,below=.15em] at (6.5,-1.75){int};
		\end{uncoverenv}
		
		\begin{uncoverenv}<4->
			\draw[thick, dashed, teal] (0,-.5) -- (3,-.5);
			\draw[thick, dashed, teal] (0,-.5) -- (0,-1);
			\draw[thick, dashed, teal] (0,-1) -- (3,-1);
			\draw[thick, dashed, teal] (3,-.5) -- (3,-1);
			\node[teal,below=.15em] at (1.5,-.5){int*};
			
			\draw (3,-.75) edge[out=0,in=180,->,shorten >=.5ex, shorten <=.5ex] (5,-1.875);
		\end{uncoverenv}
		
		\begin{uncoverenv}<5->
			\draw[thick, dashed, purple] (5,-1.25) -- (8,-1.25);
			\draw[thick, dashed, purple] (5,-1.25) -- (5,-1.75);
			\draw[thick, dashed, purple] (5,-1.75) -- (8,-1.75);
			\draw[thick, dashed, purple] (8,-1.25) -- (8,-1.75);
			\node[purple, below=.15em] at (6.5,-1.25){long};
			
			\draw[thick, dashed, teal] (0,-1) -- (3,-1);
			\draw[thick, dashed, teal] (0,-1) -- (0,-1.5);
			\draw[thick, dashed, teal] (0,-1.5) -- (3,-1.5);
			\draw[thick, dashed, teal] (3,-1) -- (3,-1.5);
			\node[teal,below=.15em] at (1.5,-1){long*};
			
			\draw (3,-1.25) edge[out=0,in=180,->,shorten >=.5ex, shorten <=.5ex] (5,-1.5);
		\end{uncoverenv}
		
		% Leave this here for historic reasons:
		%\begin{uncoverenv}<6->
		%	\draw[thick, dashed, purple] (5,-.75) -- (8,-.75);
		%	\draw[thick, dashed, purple] (5,-.75) -- (5,-1.25);
		%	\draw[thick, dashed, purple] (5,-1.25) -- (8,-1.25);
		%	\draw[thick, dashed, purple] (8,-.75) -- (8,-1.25);
		%	\node[purple, below=.15em] at (6.5,-.75){long};
		%	
		%	\draw[thick, dashed, teal] (0,-1.5) -- (3,-1.5);
		%	\draw[thick, dashed, teal] (0,-1.5) -- (0,-2);
		%	\draw[thick, dashed, teal] (0,-2) -- (3,-2);
		%	\draw[thick, dashed, teal] (3,-1.5) -- (3,-2);
		%	\node[teal,below=.15em] at (1.5,-1.5){long*};
		%	
		%	\draw (3,-1.75) edge[out=0,in=180,->,shorten >=.5ex, shorten <=.5ex] (5,-1);
		%\end{uncoverenv}
	\end{tikzpicture} \\
	\only<1>{All local variables of functions are placed at the \textit{stack}. \\
		It grows and shrinks as variables are declared and functions return.}
	\only<2>{Dynamical memory is allocated on the \textit{heap}. \\
		The example shows a function with two local \textit{int} variables.}
	\begin{onlyenv}<3>
		\begin{lstlisting}[numbers=none]
malloc(sizeof(int));
\end{lstlisting}
	Reserves exactly the amount of memory an \textit{int} variable takes.
	\end{onlyenv}
	\begin{onlyenv}<4>
		\begin{lstlisting}[numbers=none]
int *new_block = malloc(sizeof(int));
\end{lstlisting}
	The adress of that memory block is stored in an \textit{int} pointer.
	\end{onlyenv}
	\only<5>{\textit{malloc()} just needs to know the size of the block it reserves. \\
		Let us allocate a \textit{long} variable as well.}
\end{frame}

%-------------------------------------------------------------------------------------------------

\begin{frame}[fragile]{\textit{malloc()} in detail}
	The function declaration might be a little bit confusing:
	\begin{lstlisting}[numbers=none]
void *malloc(size_t size);
\end{lstlisting}
	\begin{itemize}
		\item \textit{size\_t} is an \textbf{unsigned integer} type. \\
		Any positive integer number (e.g. an \textit{int} $>0$) will do the job.
		\item \textit{size} is the size of the reserved block in \textbf{bytes}. \\
		If you want to use that block \textit{seriously}, pass the size of an actual type (e.g. \textit{sizeof(int)}).
		\item A \textit{void} pointer is returned since \textit{malloc()} does not know how you want to use the reserved block. By assigning it to a regular pointer variable it is automatically converted to that type.
	\end{itemize}
\end{frame}

%-------------------------------------------------------------------------------------------------

\begin{frame}[fragile]{Casting or not casting\dots}
	Some people claim you had to explicitly \textit{cast} the return value of \textit{malloc()}.\\
	\ \\
	\only<2->{\textbf{This is outdated.}} \ \\
	\ \\
	\begin{onlyenv}<3->
		\begin{lstlisting}[numbers=none]
int *block = malloc(sizeof(int));
\end{lstlisting}
	has the same result as
	\begin{lstlisting}[numbers=none]
int *block = (int*) malloc(sizeof(int));
\end{lstlisting}
	while the second one contains a redundant \textit{cast} and if you want to change the type of \textit{block} later, you will have to hit more keys. Consider:
	\begin{lstlisting}[numbers=none]
int *block = malloc(sizeof *block); /* gold standard */
\end{lstlisting}
	\end{onlyenv}
\end{frame}

%-------------------------------------------------------------------------------------------------

\begin{frame}{\dots that is the question}
	\huge
	\centering
	Confused? \\
	\ \\
	\only<2->{\large Do not typecast the result of malloc() \& co.}
\end{frame}

%-------------------------------------------------------------------------------------------------

\begin{frame}[fragile]{Tidying up}
	Unlike normally declared variables, dynamically allocated storage is not automatically released when the function returns.
	\begin{lstlisting}
void foo(void) {
	int *bar = malloc(sizeof *bar);
}
\end{lstlisting}
	With the pointer \textit{bar} being removed from the stack, we havo no reference on its allocated memory and those four bytes are blocked forever! \\
	\ \\
	\begin{lstlisting}[numbers=none]
free(void *ptr);
\end{lstlisting}
	Pass any pointer to previously allocated memory to \textit{free()} and it gets realeased. If you pass pointers on other things, undefined behaviour occurs (most likely program crashes).
\end{frame}

%-------------------------------------------------------------------------------------------------
\section{Dynamic arrays}
%-------------------------------------------------------------------------------------------------

\begin{frame}[fragile]{Reserving large chunks}
	To get a dynamic array of a certain \textit{type} and \textit{length}, you have to
	\begin{itemize}
		\item Pass the block size $length * sizeof(type)$ to \textit{malloc()}
		\item Assign the return value to a pointer to \textit{type}
	\end{itemize}\ \\
	\ \\
	\textit{int} array with 42 elements:
	\begin{lstlisting}[numbers=none]
int *field = malloc(42 * sizeof *field);
\end{lstlisting}\ \\
	\ \\
	Since the size of your dynamically allocated array is unknown at compile time, you cannot use \textit{sizeof} to get its length. Save it in its own variable!\\
	\ \\
	With the help of pointer arithmetic, you can use the dynamic array like a "normal" one.
\end{frame}

%-------------------------------------------------------------------------------------------------

\begin{frame}[fragile]{The fancy alternative}
		\begin{lstlisting}[numbers=none]
void *calloc(size_t nmemb, size_t size);
\end{lstlisting}
	\begin{itemize}
		\item Allocates a block of $nmemb * size$ bytes, where \textit{nmemb} is supposed to be the array's length and \textit{size} the size of its type.
		\item The whole block is filled with \textit{0}s
	\end{itemize}
	\ \\
	\begin{lstlisting}[numbers=none]
int field_length = 42;
int *field = malloc(field_length * sizeof *field);
for (int i = 0; i < field_length; i++)
	field[i] = 0;
\end{lstlisting}
	\center	$\downarrow$ Feel the difference $\downarrow$
	\begin{lstlisting}[numbers=none]
int field_length = 42;
int *field = calloc(field_length, sizeof *field);
\end{lstlisting}
\end{frame}

%-------------------------------------------------------------------------------------------------

\begin{frame}[fragile]{Resizing arrays}
	Now we come to the point that motivated us to use dynamic arrays:\\
	\ \\
	\begin{lstlisting}[numbers=none]
void *realloc(void *ptr, size_t size);
\end{lstlisting}
	\begin{itemize}
		\item \textit{ptr} is a pointer to a dynamically allocated memory block
		\item \textit{size} is the wanted new size of the memory block
		\item The return value is a pointer to the resized block
	\end{itemize} \ \\
	\ \\
	Note that the new \textit{size} can be greater or smaller than the old one!
	\begin{itemize}
		\item If it's smaller, you may lose some data at the end of the block
		\item If it's greater, the block may be at a different location in the memory $\rightarrow$ \textit{ptr} is freed then, also the additional bytes are not initialized
	\end{itemize}
\end{frame}

%-------------------------------------------------------------------------------------------------

\begin{frame}[fragile]{Clean up your code}
	Passing arrays between functions can be complicated if you store the pointer and the length seperately.\\
	\ \\
	Do you remember a way to keep different things together?\\
	\ \\
	\begin{uncoverenv}<2->
		\begin{lstlisting}[numbers=none]
struct int_array {
	int *field;
	int length;
}
\end{lstlisting}\ \\
		\ \\
		This allows you to use the \textit{struct int\_array} as a single argument or return value. Even better: pass a pointer on that structure.
	\end{uncoverenv}
\end{frame}

%-------------------------------------------------------------------------------------------------

\begin{frame}[fragile]{Strings from pointers to \textit{char}}
	By handling strings as dynamic \textit{char} arrays you can alter their size which is needed for many operations on them.
	\begin{itemize}
		\item \textit{strlen()} returns the actual length of a string (up to '\textbackslash 0' character)
		\item \textit{strncpy()} copies a string into a dynamically allocated block
	\end{itemize}\ \\
	\ \\
	\begin{lstlisting}[numbers=none]
char conststr[] = "Hello";			   /* not of much use */
int bufsize = strlen(conststr) + 1;	   /* add '\0' char */
char *str = calloc(bufsize, sizeof *str);
str = strncpy(str, conststr, bufsize); /* ready to go */
\end{lstlisting}\ \\
	\ \\
	These functions and others are declared in \textit{string.h}.
	\begin{lstlisting}[numbers=none]
$ man string.h
\end{lstlisting}
\end{frame}

%-------------------------------------------------------------------------------------------------
\section{Related Task}
%-------------------------------------------------------------------------------------------------

\begin{frame}{Vector operations}
    \href{http://fsr.github.io/c-lessons/exercises/27_vector_operations.html}{Task as online}
    \newline
    Write a program that takes two vectors as input and prints their sum.
    
    The number of elements in each vector is up to the user.
    \newline 
    \newline
    Experts: Experts: do the same with two matrices.
\end{frame}

%---------------------------------------------------------------------------------------------------------------------------
%-------------------------------------------------------------------------------------------------

\end{document}
